\documentclass{beamer}
\usetheme{metropolis}           % Use metropolis theme

\usepackage[T2A]{fontenc}
\usepackage[utf8]{inputenc}
\usepackage[english, main=russian]{babel}

\usepackage{listingsutf8}
\usepackage[ruled]{algorithm}
\usepackage{algorithmicx}
\usepackage[noend]{algpseudocode}
\usepackage{amsmath}
\usepackage{caption}
\usepackage{graphicx}
\usepackage{subcaption}
\usepackage{float}

\usepackage{tikz}
\usetikzlibrary{shapes.misc}
\usetikzlibrary{trees}

\floatname{algorithm}{Листинг}
\renewcommand{\algorithmname}{Листинг}

\title{Реализация СУБД ``ключ-значение''}
\date{\today} \author{Старовойтов Александр} \author[me]{Выполнил: Старовойтов Александр
  ИУ9-61Б\\[1mm]Руководитель: Цалкович П.А.}
% \thanks
% \institute{МГТУ им. Н.Э. Баумана}
\begin{document}
\maketitle

\begin{frame}{Цели и задачи}
  \metroset{block=fill}
	\begin{alertblock}{Цель}
      Реализовать распределенную СУБД формата ``ключ-знчение'' на основе
      структуры данных SSTables.
	\end{alertblock}
	\begin{alertblock}{Задачи}
		\begin{itemize}
          \item Изучить методы реализации распределенных СУБД;
          \item Разработать архитектуру СУБД ``ключ-значение'';
          \item Реализовать СУБД и клиент к ней;
          \item Протестировать реализацию СУБД модульными и случайными тестами.
		\end{itemize}
	\end{alertblock}
\end{frame}

\begin{frame}{NoSQL}

  \textbf{NoSQL} СУБД --- нереляционные базы данных, обычно не использующие SQL
  в качестве языка запросов.

  По модели данных их разделяют на:
  \begin{itemize}
    \item ключ-значение;
    \item столбцовые;
    \item документоориентированные;
    \item графовые.
  \end{itemize}

\end{frame}

\begin{frame}{Определение формата СУБД ``ключ-значение''}
  \textbf{СУБД ``ключ-значение''} --- это \emph{NoSQL} базы данных, с
  интерфейсом схожим с хэш-таблицей:

  \begin{itemize}
    \item \emph{get(key)}: получение данных по ключу;
    \item \emph{insert(key, value)}: вставка данных по ключу;
    \item \emph{remove(key)}: удаление данных по ключу.
  \end{itemize}

\end{frame}

\begin{frame}{Архитектура СУБД}
  \begin{alertblock}{}
    \begin{itemize}
      \item Репликация по модели ``Ведущий --- Ведомый''
      \item Глобальное версионирование всех запросов на запись
      \item Вариация LSM-дерева в качестве ядра для хранения данных
      \item SSTables для физической организации записей на диске
      \item Координация с помщью ZooKeeper
    \end{itemize}
  \end{alertblock}
\end{frame}

\begin{frame}{Модель ``Ведущий --- Ведомый''}
\begin{figure}[H]
  \centering
  \begin{minipage}[t]{\textwidth}
    \centering
    % generated by Plantuml 1.2025.2
\definecolor{plantucolor0000}{RGB}{241,241,241}
\definecolor{plantucolor0001}{RGB}{24,24,24}
\definecolor{plantucolor0002}{RGB}{0,0,0}
\begin{tikzpicture}[yscale=-1
,pstyle0/.style={color=plantucolor0001,fill=plantucolor0000,line width=0.5pt}
,pstyle2/.style={color=plantucolor0001,line width=1.0pt}
,pstyle3/.style={color=plantucolor0001,fill=plantucolor0001,line width=1.0pt}
]
\draw[pstyle0] (22.8pt,44.57pt) ellipse (8pt and 8pt);
\draw[color=plantucolor0001,line width=0.5pt] (22.8pt,52.57pt) -- (22.8pt,79.57pt)(9.8pt,60.57pt) -- (35.8pt,60.57pt)(22.8pt,79.57pt) -- (9.8pt,94.57pt)(22.8pt,79.57pt) -- (35.8pt,94.57pt);
\node at (3.5pt,96.07pt)[below right,color=black,inner sep=0]{\tiny Пользователь};
\draw[pstyle0] (124.9446pt,71.0711pt) ellipse (25.3446pt and 10.0711pt);
\node at (111.1446pt,67.0711pt)[below right,color=black,inner sep=0]{\tiny Ведущий};
\draw[pstyle0] (241.5248pt,16.0711pt) ellipse (31.2348pt and 10.0711pt);
\node at (225.5598pt,13.0711pt)[below right,color=black,inner sep=0]{\tiny Ведомый 1};
\draw[pstyle0] (241.5248pt,71.0711pt) ellipse (31.2348pt and 10.0711pt);
\node at (225.5598pt,68.0711pt)[below right,color=black,inner sep=0]{\tiny Ведомый 2};
\draw[pstyle0] (241.5248pt,126.0711pt) ellipse (31.2348pt and 10.0711pt);
\node at (225.5598pt,123.0711pt)[below right,color=black,inner sep=0]{\tiny Ведомый 3};
\draw[pstyle2] (39.89pt,71.07pt) ..controls (55.92pt,71.07pt) and (74.76pt,71.07pt) .. (93.28pt,71.07pt);
\draw[pstyle3] (99.28pt,71.07pt) -- (90.28pt,67.07pt) -- (94.28pt,71.07pt) -- (90.28pt,75.07pt) -- (99.28pt,71.07pt) -- cycle;
\draw[pstyle2] (141.88pt,63.42pt) ..controls (163.15pt,53.21pt) and (195.3106pt,37.7659pt) .. (217.4406pt,27.1459pt);
\draw[pstyle3] (222.85pt,24.55pt) -- (213.0053pt,24.8376pt) -- (218.3422pt,26.7133pt) -- (216.4666pt,32.0501pt) -- (222.85pt,24.55pt) -- cycle;
\draw[pstyle2] (150.52pt,71.07pt) ..controls (167.88pt,71.07pt) and (185.27pt,71.07pt) .. (203.85pt,71.07pt);
\draw[pstyle3] (209.85pt,71.07pt) -- (200.85pt,67.07pt) -- (204.85pt,71.07pt) -- (200.85pt,75.07pt) -- (209.85pt,71.07pt) -- cycle;
\draw[pstyle2] (141.88pt,78.72pt) ..controls (163.15pt,88.93pt) and (195.3106pt,104.3641pt) .. (217.4406pt,114.9841pt);
\draw[pstyle3] (222.85pt,117.58pt) -- (216.4666pt,110.0799pt) -- (218.3422pt,115.4167pt) -- (213.0053pt,117.2924pt) -- (222.85pt,117.58pt) -- cycle;
\end{tikzpicture}

  \end{minipage}
  \caption{Модель Ведущий --- Ведомый.}
  \label{fig:master-follwer}
\end{figure}
\end{frame}

\begin{frame}{Реализация СУБД}
  \begin{alertblock}{}
    \begin{itemize}
      \item C++
      \item Boost::asio и asio-grpc на основе корутин
      \item libzookeeper в отдельных потоках для связи с zookeeper
      \item Потокобезопасная подмена буферов на основе std::atomic
      \item Постоянное хранение с помощью ``Protocol Buffers''
      \item Вся коммуникация с клиентами и между процессами через GRPC
    \end{itemize}
  \end{alertblock}
\end{frame}

\begin{frame}{Тестирование реализации}
  \begin{alertblock}{Для проверки корректности:}
  \begin{itemize}
    \item покрытие кода юнит-тестами;
    \item случайные тесты с использованием фреймворка Jepsen.
  \end{itemize}
  \end{alertblock}
\end{frame}

\begin{frame}{Случайные тесты}
  \begin{alertblock}{}
    \begin{itemize}
      \item Реализованы в фреймворке Jepsen на Clojure
      \item Найстройка ZooKeeper и самой СУБД
      \item Кластер виртуальных машин
      \item Генератор случайных запросов
      \item Управляемые сбои
      \item Алгоритм для анализа истории выполнения операций
    \end{itemize}
  \end{alertblock}
\end{frame}

\begin{frame}{Заключение}
  \begin{alertblock}{Возможные улучшения}
  \begin{itemize}
    \item поддержка конфигурации схемы репликации;
    \item сбор метрик;
    \item переход на страничную организацию памяти с организацией кэширования;
    \item буфер для журнала предзаписи;
    \item фильтры Блума для постоянного хранилища.
  \end{itemize}
  \end{alertblock}
\end{frame}

\begin{frame}{Заключение}
  Цель курсовой работы была достигнута. Получена стабильная система,
  оптимизированная под операции записи, отвечающая требованиям доступности и
  согласованности.

\end{frame}

\begin{frame}{Заключение}
  В ходе разработки были получены следующие навыки:
    \begin{itemize}
      \item разработка многопоточных, конкурентных программ на C++;
      \item системное программирование на C++ в среде Linux;
      \item тестирование СУБД при помощи случайных тестов;
      \item реализация репликации и неизменяемых структур для постоянного
            хранения.
    \end{itemize}
\end{frame}

\end{document}
